
\documentclass{article}
\usepackage{colortbl}
\usepackage{makecell}
\usepackage{multirow}
\usepackage{supertabular}

\begin{document}

\newcounter{utterance}

\centering \large Interaction Transcript for game `wordle\_withcritic', experiment `high\_frequency\_words\_no\_clue\_no\_critic', episode 4 with Olmo{-}3.1{-}32B{-}Instruct{-}t0.0.
\vspace{24pt}

{ \footnotesize  \setcounter{utterance}{1}
\setlength{\tabcolsep}{0pt}
\begin{supertabular}{c@{$\;$}|p{.15\linewidth}@{}p{.15\linewidth}p{.15\linewidth}p{.15\linewidth}p{.15\linewidth}p{.15\linewidth}}
    \# & $\;$A & \multicolumn{4}{c}{Game Master} & $\;\:$B\\
    \hline

    \theutterance \stepcounter{utterance}  
    & & \multicolumn{4}{p{0.6\linewidth}}{
        \cellcolor[rgb]{0.9,0.9,0.9}{
            \makecell[{{p{\linewidth}}}]{
                \texttt{\tiny{[P1$\langle$GM]}}
                \texttt{You are a language wizard who likes to guess words by using the given rules.} \\
\\ 
\texttt{Welcome to Wordle! You have six attempts to guess the target word, a valid English word of five lowercase letters (a{-}z). Please use the tags "explanation:" and "guess:" to provide a concise explanation for each guess.} \\
\\ 
\texttt{For instance, if your guess is "apple", your response should be} \\
\texttt{explanation: this is a common five{-}letter English word, and I am starting my guess with this word.} \\
\texttt{guess: apple} \\
\\ 
\texttt{After each guess, your answer will be validated, and you will receive feedback indicating which letters are correct (green), which letters are correct but in the wrong position (yellow), and which letters are incorrect (red). This feedback can be useful in determining which letters to include or exclude in your next guess.} \\
\\ 
\texttt{For example, the feedback for "apple" might be:} \\
\texttt{guess\_feedback: a<yellow> p<yellow> p<green> l<yellow> e<red>} \\
\\ 
\texttt{The explanation should contain details about how the guess\_feedback is used to arrive at a new guess.} \\
\\ 
\texttt{Do not generate any other text.} \\
\\ 
\texttt{Let's begin with your first guess.} \\
\\ 
\texttt{clue: 1985 hit by heart} \\
            }
        }
    }
    & \\ \\

    \theutterance \stepcounter{utterance}  
    & \multicolumn{4}{p{0.6\linewidth}}{
        \cellcolor[rgb]{0.9,0.9,0.9}{
            \makecell[{{p{\linewidth}}}]{
                \texttt{\tiny{[P1$\rangle$GM]}}
                \texttt{explanation: The clue '1985 hit by Heart' refers to the song 'These Dreams'. However, since the target is a five{-}letter word, I will try a common five{-}letter word that includes the letter 'e' and 's', which are prominent in 'These Dreams'.} \\
\texttt{guess: these} \\
            }
        }
    }
    & & \\ \\

    \theutterance \stepcounter{utterance}  
    & & \multicolumn{4}{p{0.6\linewidth}}{
        \cellcolor[rgb]{0.9,0.9,0.9}{
            \makecell[{{p{\linewidth}}}]{
                \texttt{\tiny{[GM$\rangle$P2]}}
                \\ 
\\ 
\texttt{clue: 1985 hit by heart} \\
\texttt{explanation: The clue '1985 hit by Heart' refers to the song 'These Dreams'. However, since the target is a five{-}letter word, I will try a common five{-}letter word that includes the letter 'e' and 's', which are prominent in 'These Dreams'.} \\
\texttt{guess: these} \\
            }
        }
    }
    & \\ \\

    \theutterance \stepcounter{utterance}  
    & & & \multicolumn{4}{p{0.6\linewidth}}{
        \cellcolor[rgb]{0.9,0.9,0.9}{
            \makecell[{{p{\linewidth}}}]{
                \texttt{\tiny{[GM$\langle$P2]}}
                \texttt{Let's break down the clue and your reasoning:} \\
\\ 
\texttt{**Clue:** 1985 hit by Heart  } \\
\texttt{**Explanation:**  } \\
\texttt{The band Heart released the hit song "These Dreams" in 1985.  } \\
\texttt{However, the target is supposed to be a **five{-}letter word**.} \\
\\ 
\texttt{**Your Guess:** "these"} \\
\\ 
\texttt{But "these" is **five letters**, and it's part of the song title "These Dreams." However, "these" by itself isn't the song title, but it's a good guess given the constraints.} \\
\\ 
\texttt{But let's check if there's a more direct five{-}letter word that fits the clue, perhaps as an anagram or abbreviation, or maybe a word that's strongly associated.} \\
\\ 
\texttt{But given the clue is "1985 hit by Heart," and the only five{-}letter word that directly comes from the title is "these." Alternatively, maybe the answer is "dreams," but that's also six letters.} \\
\\ 
\texttt{Alternatively, perhaps the answer is "heart," but that's the band, not the song.} \\
\\ 
\texttt{Given the constraints, **"these"** is a reasonable guess, as it's a five{-}letter word from the song title "These Dreams," which was a 1985 hit by Heart.} \\
\\ 
\texttt{However, if the puzzle requires the actual song title in five letters, there isn't one. So "these" is the best possible answer under the five{-}letter constraint.} \\
\\ 
\texttt{**Final Answer:**  } \\
\texttt{**these**} \\
            }
        }
    }
     \\ \\

    \theutterance \stepcounter{utterance}  
    & & & \multicolumn{2}{p{0.3\linewidth}}{
        \cellcolor[rgb]{0.9,0.9,0.9}{
            \makecell[{{p{\linewidth}}}]{
                \texttt{\tiny{[GM$|$GM]}}
                \texttt{The response should always start with the keyword 'explanation:'} \\
            }
        }
    }
    & & \\ \\

    \theutterance \stepcounter{utterance}  
    & & & \multicolumn{2}{p{0.3\linewidth}}{
        \cellcolor[rgb]{0.9,0.9,0.9}{
            \makecell[{{p{\linewidth}}}]{
                \texttt{\tiny{[GM$|$GM]}}
                \texttt{game\_result = ABORT} \\
            }
        }
    }
    & & \\ \\

\end{supertabular}
}

\end{document}
